\chapter{Tonicizzazione e dominanti secondarie}
\label{chap:tonicizzazione}

Questo capitolo tratta la tonicizzazione, primo passo verso la modulazione.

\section{Modulazione e tonicizzazione: due fenomeni distinti}
\label{sec:mod_vs_ton}

\begin{definizione}[Modulazione]
    È il cambio \textbf{stabile} di centro tonale, confermato da una cadenza nella nuova tonalità. Rimanda a \ref{chap:modulazione}.
\end{definizione}

\begin{definizione}[Tonicizzazione]
    È il trattamento \textbf{momentaneo} di un grado non tonicale come fosse una tonica locale, senza abbandonare la tonalità principale. Il ritorno alla tonica originale avviene immediatamente dopo uno o due accordi.
\end{definizione}

\begin{regola}
    Se l'accordo tonicizzato è seguito da una cadenza nella nuova tonalità locale → modulazione. Se è seguito da un ritorno alla tonalità originale → tonicizzazione.
\end{regola}

\section{Il meccanismo: la dominante secondaria}
\label{sec:meccanismo_dominante}

\begin{definizione}[Dominante secondaria]
    È un accordo costruito e funzionante come dominante di un grado diverso dalla tonica. Si indica come V/II ("dominante del secondo"), V/V ("dominante del quinto"), ecc. È sempre una triade maggiore o una settima di dominante.
\end{definizione}

\begin{regola}
    Qualsiasi grado può essere tonicizzato tranne il VII (la triade diminuita non può essere tonica). I casi più comuni sono II, IV e V.
\end{regola}

\section{I casi più frequenti nel basso dato}
\label{sec:casi_frequenti}

\subsection{V/V (dominante della dominante)}
\begin{itemize}
    \item Struttura: accordo sul II grado alzato cromaticamente (\Fadiesis  in Do M).
    \item Risoluzione: su V o \sigla{V}{7}{}.
    \item Condotta: \Fadiesis  (sensibile della dominante) sale a Sol.
\end{itemize}
\musica{es_cap15_01}

\subsection{V/IV (dominante della sottodominante)}
\begin{itemize}
    \item Struttura: accordo sul I grado con settima minore (\Sibemolle  in Do M).
    \item Risoluzione: su IV.
\end{itemize}
\musica{es_cap15_02}

\section{Riconoscimento nel basso dato}
\label{sec:riconoscimento}

\begin{regola}
    Se nel basso appare un'alterazione non prevista dalla chiave, chiedersi: "questa nota è la sensibile (terza) di una dominante secondaria?". I segnali tipici sono il \Fadiesis  sul II grado (V/V) o il \Sibemolle  sul I grado (V/IV).
\end{regola}

\section{Collegamento alla modulazione}
\label{sec:collegamento_mod}

Tonicizzazione e modulazione sono due punti di un \textit{continuum}. Una tonicizzazione prolungata e confermata da una cadenza diventa a tutti gli effetti una modulazione (\ref{chap:modulazione}).
